\documentclass[a4paper,12pt]{article}
\usepackage[english]{babel}
\usepackage[utf8]{inputenc}
\usepackage[T1]{fontenc}
\usepackage{geometry}
\usepackage{hyperref}
\geometry{margin=2.5cm}

\title{\textbf{Privacy Policy on Personal Data Processing}\\
pursuant to Articles 13 and 14 of Regulation (EU) 2016/679 (GDPR)\\
NutriApp}
\date{\today}

\begin{document}

    \maketitle

    \section*{1. Data Controller}

    The Data Controller for personal data processing is \textbf{NutriApp Team}, with operational headquarters in Italy.
    For any information regarding the processing of personal data, you can contact the Data Controller at the following email address:
    \href{mailto:privacy@nutriapp.it}{privacy@nutriapp.it}.

    If a Data Protection Officer (DPO) is appointed, their contact details will be published and communicated to users.

    \section*{2. Type of Data Processed}

    As part of the use of the NutriApp application, the following categories of personal data may be processed:

    \textbf{Identification and contact data}
    name, email address, any login credentials.

    \textbf{Health-related data (special categories of data pursuant to Art. 9 GDPR)}
    weight, height, nutritional goals, eating habits, data relating to caloric and nutritional intake, any information voluntarily provided by the user regarding medical conditions or specific dietary needs.

    \textbf{Usage data}
    information relating to the use of the application, preferences, interactions with features.

    \textbf{Technical data}
    IP address, device identifiers, operating system, log data, technical cookies where applicable.

    \section*{3. Purpose of Processing}

    Personal data is processed for the following purposes:

    \begin{itemize}
        \item Providing nutritional monitoring and food diary management services;
        \item Allowing the creation and management of the user account;
        \item Processing personalized statistics and graphical visualizations of progress;
        \item Improving the quality and functionality of the application;
        \item Fulfilling legal obligations or requests from competent authorities;
        \item Managing any requests for assistance or technical support.
    \end{itemize}

    \section*{4. Legal Basis of Processing}

    The processing of personal data is based on the following legal grounds:

    \begin{itemize}
        \item Art. 6, par. 1, let. b GDPR: execution of a contract or pre-contractual measures;
        \item Art. 6, par. 1, let. c GDPR: compliance with legal obligations;
        \item Art. 6, par. 1, let. a GDPR: consent of the data subject for specific purposes;
        \item Art. 9, par. 2, let. a GDPR: explicit consent for the processing of health-related data.
    \end{itemize}

    The provision of data is necessary for the provision of the main services of the application. Failure to provide data may make it impossible to use certain features.

    \section*{5. Methods of Processing}

    Data processing is carried out using IT and telematics tools, in accordance with the principles of lawfulness, fairness, transparency, data minimization, and storage limitation.

    No automated decision-making processes are carried out that produce legal effects on the data subject pursuant to Art. 22 GDPR.

    \section*{6. Data Retention Period}

    Personal data will be kept for the time strictly necessary to achieve the purposes for which they were collected, in compliance with the principle of storage limitation referred to in Art. 5 GDPR.

    In particular:

    \begin{itemize}
        \item Data will be kept for the entire duration of the contractual relationship and the use of the application;

        \item Once the contractual relationship is terminated or the account is deleted, personal data may be kept for a further period if this is necessary to:
        \begin{itemize}
            \item fulfill legal obligations (e.g., tax or accounting obligations);
            \item protect the rights of the Data Controller, including in court;
            \item manage any complaints or disputes.
        \end{itemize}

        \item Health-related data will be kept exclusively for the time necessary to provide the service requested by the user and will be securely deleted or anonymized upon termination of the relationship, unless there is a legal obligation or other legal basis allowing their retention.
    \end{itemize}

    At the end of the periods indicated above, the data will be securely deleted or irreversibly anonymized.

    \section*{7. Data Communication and Transfer}

    Personal data may be communicated to:

    \begin{itemize}
        \item IT and hosting service providers;
        \item Cloud infrastructure providers;
        \item Legal or tax consultants, where necessary.
    \end{itemize}

    These entities operate as Data Processors pursuant to Art. 28 GDPR.

    If data is transferred to third countries outside the European Economic Area, the transfer will take place in compliance with Articles 44 and following of the GDPR, by adopting appropriate safeguards (e.g., Standard Contractual Clauses).

    \section*{8. Rights of the Data Subject}

    The data subject may exercise the rights provided by Articles 15-22 GDPR at any time, including:

    \begin{itemize}
        \item Right of access to personal data;
        \item Right to rectification;
        \item Right to erasure (right to be forgotten);
        \item Right to restriction of processing;
        \item Right to data portability;
        \item Right to object to processing;
        \item Right to withdraw consent at any time.
    \end{itemize}

    Requests can be sent to the email address \href{mailto:privacy@nutriapp.it}{privacy@nutriapp.it}.

    It is also possible to lodge a complaint with the Data Protection Authority.

    \section*{9. Data Security}

    NutriApp adopts appropriate technical and organizational measures pursuant to Art. 32 GDPR, including:

    \begin{itemize}
        \item Secure user authentication;
        \item Data access limited to authorized personnel;
        \item Periodic backups and protection systems against unauthorized access.
    \end{itemize}

    \section*{10. Minors}

    The application is not intended for children under 14 years of age.

    Pursuant to Art. 8 of the GDPR and applicable local laws, the processing of personal data based on consent in relation to the direct offer of information society services is lawful only if the child is at least 14 years old.

    If the user is under 14 years of age, the processing of personal data is lawful only if consent is given or authorized by the holder of parental responsibility.

    The Data Controller may adopt reasonable measures to verify the age declared by the user and, where necessary, the existence of consent from the parent or guardian.

    If it is ascertained that personal data of a minor under 14 has been collected without the required consent, such data will be promptly deleted.

    \sectionsection*{11. Changes to this Privacy Policy}

    This privacy policy may be updated over time.
    Any substantial changes will be communicated to users via in-app notification or email.

    \vspace{1cm}

    \noindent
    Last updated: \today

\end{document}
